\section{The Scaling Numbers as a Bayes Net Inventory}
\label{sec:scaling-numbers}

The five architectural dimensions of a transformer --- $W$ (sequence length),
$D_{\mathrm{model}}$ (embedding dimension), $H$ (attention heads), $D_{ff}$ (FFN
hidden dimension), $L$ (layers) --- are not arbitrary engineering choices. In the
Bayes net view, each dimension has a precise interpretation, and together they
constitute an inventory of the implicit factor graph.

\begin{center}
\begin{tabular}{lll}
\toprule
Dimension & BP interpretation & Llama~3 405B \\
\midrule
$W$ & nodes in the factor graph (sequence positions) & 8{,}192 \\
$L$ & rounds of belief propagation & 126 \\
$H$ & AND patterns per layer & 128 \\
$D_{ff}$ & OR patterns per layer & 53{,}248 \\
$L \times H$ & distinct AND weight patterns & 16{,}128 \\
$L \times D_{ff}$ & distinct OR weight patterns & 6.7M \\
$L \times H \times W$ & AND evaluations per forward pass & 132M \\
$L \times D_{ff} \times W$ & OR evaluations per forward pass & 55B \\
\bottomrule
\end{tabular}
\end{center}

The ratio $D_{ff} / H \approx 416$ for Llama~3 405B. The architecture devotes
$416\times$ more capacity to inference (OR gates) than to routing (AND gates).
This is the right ratio for a system whose bottleneck is combining evidence, not
finding it.

\subsection*{What Scaling Adds}

When a model scales from 8B to 405B parameters, what changes in Bayes net terms?

\begin{itemize}
\item \textbf{More rounds of BP} ($L$ increases): deeper reasoning chains become
possible. A 126-layer model can follow implications 126 hops deep.

\item \textbf{Richer OR patterns} ($D_{ff}$ increases): more conditional probability
table entries per layer. The model can represent more fine-grained conditional
relationships.

\item \textbf{More nodes} ($W$ increases): longer sequences mean larger factor
graphs. The model reasons over more propositions simultaneously.

\item \textbf{AND patterns stay small} ($H$ grows slowly): 128 heads in 405B vs.\
16 in the smallest models. The number of distinct routing operations is nearly
constant. What scales is inference capacity, not routing capacity.
\end{itemize}

The implication is sharp: scaling adds inference depth and richness, not
fundamentally new reasoning structure. The AND/OR architecture is fixed. What
changes is how many rounds it runs and how many patterns it can match.

\subsection*{The 16,128 Number}

The 16,128 distinct AND weight patterns for Llama~3 405B is a striking number
because it is so small. Every token position in every sequence uses one of these
16,128 routing patterns --- the same patterns, regardless of what the token is.
The routing key (which neighbor to attend to) is determined by the weight matrices,
which are shared across positions. What varies across positions is the continuous
content of the residual stream, not the routing structure.

This is the formal statement of weight sharing as universal rule implementation
(Section~4.3 of~\cite{coppola2026transformers}): the same computation runs for every
instantiation of a concept, and the particular input supplies the magnitudes. The
16,128 patterns are the universal rules. The residual stream values are the
particular instances.